\newpage
\section{Exact calculations in the spherical module}
\label{sec:checks}

The accompanying Python program implements two calculations over
$\mathbb Z[v,v^{-1}]$. In the first, represent a standard vector by the
levels $a_i=\lfloor(x(i)-1)/N\rfloor$ of the left-minimal window.
To apply $C_P/\pi_T$, where $T$ is the previous intersection, cut the
circle at a missing reflection and add one to the levels that pass the cut.
In each position block sort the lifted levels in increasing order. Let $L$
be the stabilizer of this minimum in $W_P$, $d(x)$ the within-block
inversion number and $D=\ell(w_P)-\ell(w_L)$. The new coefficient at $y$
in its orbit $O$ is
\[
 v^{D-d(y)}\frac{\pi_L}{\pi_T}\sum_{x\in O}v^{-d(x)}[M_x]F .
\]
Indeed $M_xC_P=v^{-d(x)}M_{\min O}C_P$ and splitting $W_P$ into
minimal left $W_L$-coset representatives evaluates the latter expression.
All Laurent divisions in the program check for zero remainder.

The second calculation uses only right simple reflections. For a right
descent $s$ of $x$, recursively calculate $\mathbf C_{xs}C_s$. At
each strict lower $z$ subtract the constant term of its standard
coordinate times $\mathbf C_z$. The resulting monic bar-invariant
vector satisfies the canonical triangularity condition. This recursion independently computes the seed in the parabolic paper, every endpoint, and the canonical transition matrix
used to decompose the products. The numerical inputs are the parabolic
sets, the chain words and the endpoint windows.

Let $L_m=C_{s_2}\cdots C_{s_m}C_{s_1}\cdots C_{s_{m-1}}$ and let
$D_m$ denote the two normalized parabolic averages. Each cell below
gives the degree-one identity coefficient followed by the number of
distinct canonical indices with nonzero graded multiplicity.
\begin{center}
\begin{tabular}{c|ccccc}
$m$ & $\mathcal R_m$ & $\mathcal{BS}_m L_m$ &
$\mathbf C_{X_0}L_m$ & $\mathbf C_{X_0}D_m$ & $\mathbf C_{\beta_m}$\\ \hline
2 & $4,5$ & $4,5$ & $3,3$ & $3,3$ & $2,1$\\
3 & $11,11$ & $18,23$ & $12,13$ & $7,6$ & $3,1$\\
4 & $25,17$ & $81,70$ & $50,38$ & $14,9$ & $4,1$
\end{tabular}
\end{center}
In each rank all five columns have identical standard support,
respectively $38,368,3098$ indices, and coefficient one at $M_{\beta_m}$.
Every coordinate of $\mathcal{BS}_mD_m$ agrees with $\mathcal R_m$.
Of the family-F summands, the first two columns contain the entire row
with multiplicity one, and the diagonal is the unique family-F term in the last three.
The independent KL recursion gives rows $(1,2),(1,2,3),(1,2,3,4)$.
For instance
\[
 [M_e]\mathbf C_{\beta_2}=2v+4v^3+5v^5+4v^7+v^9,\qquad
 [M_e]\mathbf C_{X_0}=v+v^3+2v^5+v^7\quad(m=2).
\]
Executing \texttt{python3 compare.py --full 2 3 4} in the directory
containing \texttt{compare.py} reproduces the tables and full identity
polynomials in \texttt{results.txt} without external packages. A second
script, \texttt{python3 partial.py 2 3 4 5}, checks the mixed
full-spherical coordinates and the discarded degrees directly. In rank
five its successive corrected degree-one coefficients are $1,3,6,10,15$;
the independent KL recursion gives $(1,2,3,4,5)$. For selected entries at larger values of $m$, a separate run gives $\mu_{2;6}=\mu_{2;7}=\mu_{2;8}=2$, $\mu_{3;6}=\mu_{3;7}=3$ and $\mu_{4;6}=4$. The last expansion has 120\,862 standard terms. The script \texttt{check\_off\_diagonal.py} and its output accompany the calculation.

\paragraph{Comparison after boundary projection.}
At $m=2$, $L_m=D_m$. For $m\ge3$, applying the projector $C_V/\pi_V$ to
$\mathcal{BS}_mL_m$ gives a vector linearly independent of
$\mathcal R_m$ over $\mathbb Q(v)$. It suffices to compare two
coordinates at $v=1$. Write $t=2^{m-1}/m$. The ratio at $M_e$ is
$t^2$ since each simple factor multiplies this coordinate by 2,
whereas each divided parabolic factor multiplies it by $m$.
The ratio at $M_{\beta_m}$ is $t$, as follows.

In type $m$ the first affine product has just the word $A^mBC^m$.
The next $K_m$-average fixes the initial $A$ and puts the $B$ in any
one of positions $2,\ldots,m+1$. The $E$-average then gives coefficient
one to every type-$m$ word with prefix $AB$ (the stabilizer of the
remaining letters has size $(m-1)!m!$, exactly the overlap denominator).
The last $m$ letters stay fixed during $L_m$. In the $V$-orbit of
$\beta_m$ they contain one $C$ and $m-1$ letters $A$, so only
the initial prefix $ABC^{m-1}$ of length $m+1$ can contribute:
the leftmost letter is $A$ and the second letter cannot be $C$.
All first-chain swaps must be chosen to put $B$ at position $m+1$.
Each of the $2^{m-1}$ choices in the second chain then contributes.
There are $m$ possible suffixes and $m^2$ words in the target orbit.
The normalized $V$-average gives $m2^{m-1}/m^2=t$ at $M_{\beta_m}$,
where $\mathcal R_m$ has coefficient one. The denominators and the
coordinate used to compare scalars are nonzero at $1$. Since $t>1$,
the two ratios differ.

\paragraph{Partial galleries.}
For $T_j=\mathcal{BS}_jD_m$ with $1\le j\le m\le4$, exact
elimination gives family-F multiplicity one exactly at indices
$1,\ldots,j$. A truncated gallery can also contain canonical indices
outside the six families for $\mathcal R_m$. At $(m,j)=(3,2)$ its
top word is $CCABBBA$, and the family-F word $CBABCBA$ occurs in a lower
summand. Section~\ref{sec:partialrows} singles out the positive summand
supported on the relevant seed indices in every rank.


\subsection*{Affine-window implementation}
The directory \path{comparison/independent_checks/} contains an exact
affine-window implementation separate from the level-word implementation
above. Its inputs are the defining Hecke relations, parabolic subsets and
endpoint windows. The recursion constructs canonical vectors directly
from these inputs, giving an independent check on the row formula and
the canonical expansions.

The engine's validation checks quadratic, braid and commuting relations,
inversion, and parabolic-quotient enumeration.
The normalized partial-sum calculation at \(m=2,3,4,5\) checks every
standard coordinate of \(\mathcal{BS}_mT_1T_2=\mathcal R_m\), every
family-F multiplicity in every \(Y_j\), every spherical correction
\(g_{j-r}\), and the degree bounds on the remaining terms.
The corrected partial rows are \(1,3\), \(1,3,6\), \(1,3,6,10\),
and \(1,3,6,10,15\).
The finite-tail calculation checks the canonical identities in
Section~\ref{sec:finite-tails} through \(m=5\), including all lower
correction terms.

The scripts and their saved outputs make these exact Python calculations
reproducible. They check finite instances of the uniform arguments given
in the note and are separate from our Lean development.
